\section{Benchmark Design}
\label{sec:benchmark}

\subsection{Problem Types}

ECAG-Bench problems fall into four categories:

\begin{description}[nosep]
  \item[Example:] Construct a mathematical object with specified properties. (``Give an example of a flat morphism that is not smooth.'')
  \item[Counterexample:] Construct an object that violates a plausible-sounding statement. (``Give a counterexample to: every coherent sheaf on a smooth variety is locally free.'')
  \item[Computation:] Compute a specific invariant of a given object. (``Compute the Picard group of $\Spec \mathbb{Z}[\sqrt{-5}]$.'')
  \item[Proof:] Prove a statement about a specific object. (``Show that the Horrocks--Mumford bundle is stable.'')
\end{description}

\subsection{Four-Layer Verification Pipeline}
\label{sec:verification}

Since constructions in algebraic geometry are non-unique, standard answer-matching is insufficient.
ECAG-Bench employs a four-layer verification pipeline:

\paragraph{Layer 1: Parsing.}
The model's output is parsed to extract the claimed mathematical object.
This layer uses regex and SymPy to determine whether the output is a well-formed mathematical expression (a polynomial, a variety defined by equations, a morphism given by coordinate functions, etc.).

\paragraph{Layer 2: CAS Verification.}
The extracted object is verified using computer algebra systems---primarily SageMath, with Singular and Macaulay2 for specialized computations.
The CAS checks whether the constructed object actually satisfies the required properties: genus, smoothness, dimension, cohomology vanishing, etc.
This is the \emph{core layer}: 309 of 332 problems (93.1\%) use CAS verification.

\paragraph{Layer 3: LLM-Assisted Judgment.}
For problems where CAS verification is insufficient (e.g., those requiring assessment of reasoning quality or non-triviality), a rubric-based LLM judge evaluates:
\begin{itemize}[nosep]
  \item (C1) Correctness: Is the constructed object valid?
  \item (C2) Completeness: Are all required properties verified?
  \item (C3) Non-triviality: Is the answer non-degenerate?
  \item (C4) Reasoning: Is the justification sound?
\end{itemize}

\paragraph{Layer 4: Human Expert Sampling.}
A random subset of model responses is reviewed by a human expert in algebraic geometry to calibrate the automated layers and catch systematic errors.

\subsection{Verification Strategies}

Each problem is assigned one of three verification strategies:

\begin{description}[nosep]
  \item[Strategy A (Exact):] Problems with unique numeric or symbolic answers (e.g., ``Compute $\dim H^1(X, \mathcal{O}_X)$'').
  Verified by exact match. 14 problems (4.2\%).
  \item[Strategy B (CAS):] The model outputs a mathematical object; CAS verifies its properties deterministically.
  This is the primary strategy. 309 problems (93.1\%).
  \item[Strategy C (Rubric):] Structured scoring on correctness, completeness, non-triviality, and reasoning validity.
  Used for proofs and complex constructions. 9 problems (2.7\%).
\end{description}

\subsection{Evaluation Protocol}

Models are evaluated in a zero-shot, single-turn setting.
Each problem is presented as a standalone prompt with the following format:

\begin{quote}
\texttt{You are an expert in algebraic geometry. [Problem statement]. Provide a concrete construction with full justification.}
\end{quote}

Responses are scored as:
\begin{itemize}[nosep]
  \item \textbf{Correct (1.0):} The constructed object is valid and satisfies all stated properties.
  \item \textbf{Partial (0.5):} The object is valid but missing some properties, or the justification is incomplete.
  \item \textbf{Incorrect (0.0):} The object is invalid, does not exist, or does not satisfy the required properties.
\end{itemize}
